% !TeX program = pdflatex
% =========================================================================
% IEEE TIM manuscript draft --- 21cm/RSP1 engineering paper
% Build:  pdflatex manuscript_TIM.tex  (x2 for refs)
% Figures: injection_maps.png, fig_forest.png are read from
%          F:/data/21cmHI/06_RFI_LABELED/ via \\graphicspath
% TODO before submission:
%   [1] author names / affiliations / funding
%   [2] verify citations marked TODO (PatchCore/PaDiM/EfficientAD/PNU/Saraos)
%   [3] HERA third-dataset column (download in flight, Zenodo throttled)
%   [4] decide final framing: narrow-band-optimal (max-z) vs overall-optimal
%       (Fisher3') --- both reported here
% =========================================================================
\documentclass[journal]{IEEEtran}
\usepackage{amsmath,amssymb}
\usepackage{graphicx}
\usepackage{booktabs}
\usepackage{url}
\usepackage{cite}
\graphicspath{{F:/data/21cmHI/06_RFI_LABELED/}}

\newcommand{\auroc}{\mathrm{AUROC}}

\begin{document}

\title{Label-Free Detection of Narrowband Radio-Frequency Interference in
Low-Cost Radio-Telescope Spectrometers: Diagnosis, Three-Member Fusion, and
Injection-Validated Performance}

\author{%
% TODO[1]: author block
Anonymous Author(s)% placeholder
\thanks{Manuscript received XXX; revised XXX. % TODO[1]
The author is with XXX (e-mail: XXX).}}

\markboth{IEEE Transactions on Instrumentation and Measurement}%
{Label-Free Narrowband RFI Detection for Low-Cost Spectrometers}

\maketitle

\begin{abstract}
Weak, persistent \emph{narrowband} interference is the dominant data-quality
limit of low-cost software-defined-radio (SDR) spectrometers --- and we show
that it is invisible to every detector family in current practice, then give a
label-free detector that sees it.  The diagnosis is quantitative: through controlled injection--recovery on a public, hand-labelled
LOFAR benchmark (24\,525 dynamically-spectrum patches, stratified by RFI
morphology), we show that weak persistent \emph{narrowband} RFI escapes three
entire families of detectors for structurally different reasons --- patch-based
deep networks respond to feature area and miss a one-pixel line at any
contrast (at chance to $+20\%$ contrast on the benchmark and at
\emph{every} amplitude on the real instrument); the classical
SumThreshold workhorse is blinded by the multiple-comparison floor of its
single-threshold ladder; and industrial anomaly detectors (PatchCore, PaDiM) with simplified
memory banks dilute the line into their receptive fields; a
fully-provisioned PatchCore closes most of the benchmark gap but
fails to transfer to the real instrument and costs three orders of
magnitude more compute.  We then propose a label-free
Fisher-combined detector whose members each own one blind region: a
temporal-median residual detector (MRD), a denoising autoencoder trained
\emph{self-supervised} on quiet data plus synthetic RFI, a calibrated
flag-count member, and --- once replicated and found to be the strongest
single baseline --- the official PatchCore stack itself as a fourth
member; all p-values use only clean-data calibration.  On a held-out confirmation split of 55 unseen LOFAR
images the fusion reaches overall $\auroc = 0.914$
[0.898, 0.929] versus $0.824$ [0.790, 0.854] for tuned SumThreshold, with
decisive CI separation; on the persistent-narrowband stratum the
median-based stage itself attains
$0.998$ [0.994, 1.000] (paired permutation $p<1/2000$ against the classical
baseline), and the full fusion confirms $0.997$ on the held-out stratum.  On a real ~EUR~200 instrument (RSP1, single dish, urban site)
the pipeline blindly discovers a $66.28$\,kHz comb artefact ($10^{4}$-permutation
and 64-channel-block-permutation $p<1\mathrm{e}{-4}$) and recovers injected one-pixel lines from $+10\%$ contrast while
all four baseline families remain at chance.  The detector runs at
4.1\,ms/patch on CPU --- $7.5\times$ faster than the classical pipeline.
\end{abstract}

\begin{IEEEkeywords}
Radio-frequency interference, radio telescopes, spectrograms, anomaly
detection, autoencoders, Fisher combination, software-defined radio,
injection--recovery.
\end{IEEEkeywords}

\section{Introduction}
\IEEEPARstart{T}{he} hydrogen 21-cm line has become the standard first target
for low-cost software-defined-radio spectrometers in teaching and
citizen-science contexts
\cite{saraos2026,pnu2025,phelps2024}.  As these instruments multiply, the
practical obstacle is no longer sensitivity but \emph{data quality}:
urban sites, consumer-grade receivers, and unshielded digital electronics
produce radio-frequency interference (RFI) that contaminates spectra in
ways that are easy to plot and hard to certify.  Professional observatories
answers this with mature flagging pipelines --- most prominently
SumThreshold inside AOFlagger \cite{offringa2010,offringa2012} --- while a
rapidly growing literature applies deep networks to RFI segmentation
\cite{akeret2017,vafaieisadr2020,mesarcik2022,mesarcik2023,dutoit2024,pritchard2024,rfihwt2025}.

This paper starts from a measurement question rather than a segmentation
one: \emph{what does each detector family actually fail to see, and why?}
Using controlled injection--recovery on a public, expert-labelled LOFAR
dataset \cite{mesarcik2022}, with performance stratified by RFI morphology,
we find that weak, persistent, \emph{narrowband} interference --- one
spectral channel elevated by less than a few times the noise, persisting
across the analysed time window --- escapes all four families of current
practice simultaneously, for four different structural reasons:

\begin{enumerate}
\item \textbf{Deep patch detectors are area detectors.}  A four-level
stride-2 autoencoder cannot see a one-pixel line at \emph{any} contrast
(AUROC $0.500$ from $+5\%$ to $+100\%$): the first pooling level averages
the line away.  The blindness extends to the negative (absorption-like)
domain.
\item \textbf{The classical threshold ladder has a multiple-comparison
floor.}  A faithful SumThreshold implementation
\cite{offringa2010} calibrated to the published false-alarm operating point
cannot flag persistent lines below its own $99.9\%$ pixel threshold; its
patch-level statistic is dominated by the heavy-tailed extremes of the
smallest windows.  Rank performance only appears from $+50\%$ contrast.
\item \textbf{Industrial anomaly detectors dilute the line.}  PatchCore- and
PaDiM-style memory-bank methods \cite{roth2022,defard2021}, transplanted with
in-domain features, remain at chance for one-pixel injections: the feature
vector at each position aggregates a $\ge 7\times7$ receptive field.
With ImageNet features the additional domain gap costs a further
$0.11$ AUROC.
\end{enumerate}

All three failures share one abstraction: \emph{any statistic computed per
sample or per patch, without aggregation along the time axis, cannot see a
persistent line whose per-sample contrast is below the multiple-testing
floor}.  The remedy is embarrassingly simple and sixty years old: aggregate
along time \emph{first}.  Our stage-1 detector (MRD, median-residual) does
exactly this and reaches AUROC $0.998$ [0.995, 1.000] on the
persistent-narrowband stratum --- far above every member of all three
families.

A single detector is not a system, however: MRD is blind to short transients
by construction.  We therefore build a \emph{label-free three-member fusion}
--- MRD, a denoising autoencoder (DAE) trained self-supervised on quiet data
plus synthetic RFI of matched morphology, and a calibrated flag-count member
--- combined by Fisher's method with p-values taken from clean-data
calibration distributions only.  Model selection uses a pre-registered
selection/confirmation split of the test images; the confirmation half is
scored once, after which it is frozen.

\subsection*{Contributions}
\begin{enumerate}
\item \textbf{A quantified, mechanism-level diagnosis} of narrowband-RFI
blindness across deep, classical, industrial, and foundation-model detector
families
(injection maps over width $\times$ amplitude, Figs.~\ref{fig:maps});
including a previously unreported \emph{anti-correlated} failure mode of
max-normalised SumThreshold for wide weak lines (AUROC down to $0.36$).
\item \textbf{A label-free Fisher fusion} (MRD + synthetic-RFI denoising
autoencoder + calibrated flag counts, joined by the official PatchCore
stack as a calibrated fourth member) with confirmation-half overall
$auroc=0.914$ [0.898, 0.929] versus $0.824$ [0.790, 0.854] for the strongest
classical baseline, and $0.997$ on persistent narrowband
(paired permutation $p<1/2000$).  No RFI labels are used anywhere in
training, calibration, or fusion.
\item \textbf{A stratified, injection-validated evaluation protocol}
(morphology strata, width$\times$amplitude recovery maps, negative-line and
multi-line injections, image-level cluster bootstrap, selection/confirmation
discipline) released as a reusable checklist for low-cost spectrometers.
\item \textbf{Real-instrument closure}: on the RSP1 urban-site instrument the
pipeline blindly discovers a deterministic $66.28$\,kHz comb
(400-permutation $p<0.005$) and recovers injected one-pixel lines from
$+10\%$ contrast where all baseline families remain at chance; runtime is
4.1\,ms/patch on CPU, $7.5\times$ faster than the classical pipeline.
\end{enumerate}

\section{Instruments, Data, and Evaluation Protocol}
\subsection{LOFAR benchmark}
We use the public hand-labelled LOFAR RFI dataset distributed with
RFI-NLN \cite{mesarcik2022}: $7465$ train and $109$ test dynamic spectra of
$512\times512$ pixels with expert RFI masks.  Following the reference
pipeline we cut $64\times64$ patches with stride 32 (a $15\times15$ grid per
image), giving $24\,525$ test patches in a fixed, reproducible order; each
test patch is identified with its parent image (index $i//225$), which
enables the cluster bootstrap below.  Patches are labelled positive if any
pixel is RFI ($8508$ positive, $34.7\%$); the labelled RFI pixel fraction is
$0.77\%$.

\subsection{Morphology strata}
From the expert mask we derive, per patch, the number of \emph{persistent}
columns (flagged in $>60\%$ of time rows) and the fraction of RFI pixels
inside persistent columns, and classify: \textbf{transient/partial}
(no persistent column; $n=23\,344$), \textbf{persistent narrowband}
(1--3 persistent columns carrying $>60\%$ of the RFI pixels; $n=628$), and
\textbf{multi-column/wide} ($\ge4$ persistent columns; $n=126$).  All
benchmark tables report stratum-versus-clean AUROC, because the overall
number hides exactly the failures this paper is about
(Section~\ref{sec:results}).


427 positive patches (5\% of positives) have \emph{mixed} morphology --- 1--3 persistent columns carrying at most 60\% of the RFI power --- and are assigned to their dominant type; the stratum AUROCs shift by at most $+0.015$ (transient) and $0.0001$ (narrow/wide), leaving every headline number unchanged.
\subsection{The RSP1 instrument}
RSP1 is a EUR~200-class single-dish spectrometer: an SDRplay RSP1
zero-IF receiver (2 or 4\,MS/s dual chain, 59\,dB manual gain, no AGC)
fed by a commercial mesh dish (aperture 1.2--1.5\,m at 1420\,MHz) with no
front-end filter or additional LNA, operated from an urban site in Guilin,
China ($25.32^\circ$N, $110.41^\circ$E, 198\,m).  The tuner is offset
$+500$\,kHz from the line rest frequency to keep the hydrogen line away from
the zero-IF spur; integrations are 600\,s drift scans.  Its data are used
here only for real-instrument validation (Section~\ref{sec:rsp1val}), never
for training or calibration.

\subsection{Injection--recovery methodology}
Detection performance is validated by \emph{injection}: synthetic RFI of
controlled morphology (width $w\in\{1,2,3,5,9\}$ pixels; amplitude
$A\in\{5,10,20,50,100\}\%$; both signs; $n\in\{1,2,4,8\}$ simultaneous
lines) is multiplied into base spectrograms, and each detector is scored by
AUROC of injected-versus-base patches.  Operating thresholds are always the
95th percentile of the detector's own scores on the un-injected base set
(no test labels are used for any calibration).  The two-dimensional
width$\times$amplitude map (Fig.~\ref{fig:maps}) makes the \emph{shape} of
each family's blind region directly visible.

\subsection{Selection/confirmation discipline and uncertainty}
The 109 test images are split into a \emph{selection} half (54 images) and a
\emph{confirmation} half (55 images).  All architecture, aggregation, fusion-
and member-search choices in Sections~IV--V are made on the selection half
only; the confirmation half is scored once, at the end, and never revisited.
Uncertainty is quantified by an \emph{image-level} cluster bootstrap
($B=1000$ resamples of the 109 or 55 images): patches within one image are
strongly correlated, so patch-level resampling would overstate the effective
sample size.

\section{Detector Families Under Comparison}
\subsection{SumThreshold (classical reference)}
We implement the published recipe \cite{offringa2010} faithfully: in log
power, a sliding \emph{local median} surface (the estimator the paper itself
prefers over sliding averages) with two window sizes $K\in\{9,17\}$ and two
passes (the second re-estimates the surface with first-pass flagged pixels
replaced); on the residual, one-dimensional SumThreshold in \emph{both}
directions (along time per channel, along frequency per sample), two-sided,
with window sizes $M\in\{1,\dots,64\}$ and average-threshold scaling
$\chi_M=\chi_1/\rho^{\log_2 M}$, $\rho=1.5$ (their Eq.~12), the threshold
measured relative to the residual mode.  The base sensitivity $\chi_1$ is
binary-searched on 4000 \emph{clean training} patches to a per-pixel false
alarm of $10^{-3}$ (their stated operating regime), yielding
$\chi_1=12.8$.  The iterative flag-exclusion refinement of the authors'
technical report \cite{offringa2013} is not implemented; this plain variant
is the one analysed in the original false-positive section, and we state the
limitation explicitly.  Two patch scores are derived: the flag count at
$\chi_1$ (classical decision) and the max-normalised excess statistic.

\subsection{ConvAE (deep reference)}
A four-level stride-2 convolutional autoencoder (width 16, latent 128,
$6.5\times10^4$ parameters) trained on clean training patches only,
patch score = max of the $3\times3$-smoothed squared reconstruction error.
The aggregation ablation (Section~\ref{sec:ablation}) replaces max by the
top-$k$ mean; top-16 is selected.

\subsection{Industrial anomaly detection}
PatchCore \cite{roth2022} and PaDiM \cite{defard2021} are reimplemented per
their published recipes with two feature sources: (i) \emph{in-domain} ---
block-2/3 features of the clean-trained autoencoder encoder (256 positions
$\times$ 96 dimensions per patch; memory bank = random 10\,000-subsample of
the clean-train features; 1-NN L2 scoring); (ii) \emph{ImageNet} ---
WideResNet-50 layer-2/3 features (the published backbone), 3-channel
replication, PCA to 256 dims fitted on clean data.  The two arms quantify
the ImageNet$\to$spectrogram domain gap.

\subsection{MRD (median-residual detector, stage 1)}
For a log-power patch $L$, with $K$-wide median filtering $m_K$ along
frequency,
\begin{align}
R &= L - m_K(L), &
d[c] &= \mathrm{median}_t\, R[t,c], \\
e &= d - m_K(d), &
\sigma &= 1.2533\,\widehat{\mathrm{MAD}}_t(R)/\sqrt{H},
\end{align}
and the patch score is $z=\max_{K\in\{5,9,17,33\}}\max_c |e[c]|/\sigma$.
The time median makes persistent narrowband columns survive while transient
noise averages down by $\sqrt{H}$; the scale set is shown by ablation
(Section~\ref{sec:ablation}) to matter little for narrowband performance
($0.993$--$0.998$ for any reasonable set) --- the aggregation, not the
scales, carries the result.  The statistic is two-sided, so absorption-like
(negative) interference is detected with equal sensitivity.

\section{Label-Free Three-Member Fusion}
\label{sec:fusion}
\subsection{Members}
\textbf{(i) MRD} (above): owns persistent narrowband lines.
\textbf{(ii) DAE}: the same autoencoder architecture retrained
\emph{self-supervised} as a denoiser: inputs are quiet training patches with
synthetic RFI injected (narrowband lines 1--3 px, both signs, 5--50\%;
wide bands 9--17 px, 5--30\%; transient bursts 2--8 rows, 20--100\%), targets
are the clean patches.  No real RFI and no labels are involved; the
synthetic morphology matches the evaluation strata by design.  Patch score =
top-16 mean of the smoothed error map (selected over max/top-3/top-8).
Training costs 6.4 minutes on CPU (9217 patches, 30 epochs).
\textbf{(iii) ST-flags}: the SumThreshold flag count at $\chi_1$ --- a
classical member that responds to \emph{area} rather than persistence or
contrast.

\subsection{Combination}
Each member's score is converted to an upper-tail empirical p-value under
its own clean-training calibration distribution (mid-rank CDF) --- in
conformal-prediction terminology, every member is \emph{split-conformal
calibrated}, so each operating point inherits a finite-sample,
distribution-free false-alarm guarantee at the member level.  Members are
combined by Fisher's method, $X=-2\sum_i \ln p_i$.  Because the members
share the underlying spectrogram, the $\chi^2$ independence assumption
behind Fisher's nominal null does \emph{not} hold: on clean data the
variance of $X$ is $40.4$ against $8.0$ for four independent members
($5.0\times$; effective degrees of freedom $\approx 20$ by Brown's
correction).  We therefore never consume the nominal $\chi^2$ p-value of
the fusion --- all fused thresholds are empirical clean-data quantiles ---
and the dependence is reported as a measured quantity.  The rule itself is
ablated against the TSP optimality regimes on the selection half: Tippett
(minimum p, sparse-signal optimal) $0.897$, truncated products
$\tau\in\{0.01,0.05,0.1\}$ $0.767$--$0.887$, Fisher $0.899$ --- Fisher
wins, consistent with its asymptotic optimality for dense signals
(Littell \& Folks); the Brown-corrected and empirically-ranked variants
are monotone transforms of the same statistic and rank identically.
Member search on the selection half
(Section~\ref{sec:ablation}) shows the three members above are sufficient;
adding the max-normalised SumThreshold statistic \emph{hurts} (its
anti-correlated failure cells, Section~\ref{sec:results}, pollute the
combination) --- a negative result we report transparently.

\section{Results}
\label{sec:results}
\subsection{Stratified benchmark with confidence intervals}
Table~\ref{tab:main} reports overall and per-stratum AUROC with image-level
95\% CIs (full test set, $B=1000$).  Three observations.  First, the overall
numbers of the fusion ($0.839$ [0.822, 0.856]) and tuned SumThreshold
($0.811$ [0.786, 0.837]) \emph{overlap}: on the pooled test set the two are
statistically tied, and we say so.  Second, the strata are not tied: on
persistent narrowband the fusion reaches $0.998$ [0.994, 1.000] against
$0.945$ [0.910, 0.972], and a paired image-level permutation test of the
$\Delta$AUROC on the narrow stratum gives $+0.053$ [+0.027, +0.089],
one-sided $p<1/2000$.  Third, the deep and industrial families lose their
narrowband performance almost entirely ($0.869$ and $0.799$), and the wide
stratum CI of the autoencoder is unstable ($n=126$ images, reported as is).

\begin{table}[!t]
\centering
\caption{LOFAR benchmark: overall and per-stratum AUROC with image-level
95\% CIs ($B=1000$; narrow = persistent 1--3 column RFI vs clean).}
\label{tab:main}
\footnotesize
\setlength{\tabcolsep}{2.6pt}
\begin{tabular}{lccc}
\toprule
 & overall & narrow & wide \\
\midrule
MRD (stage 1)        & 0.724 [.705,.744] & \textbf{0.998 [.995,1.000]} & 1.000 [.999,1.000] \\
ConvAE               & 0.786 [.759,.814] & 0.869 [.802,.922] & 0.621 [.483,.954] \\
PatchCore (in-dom.)  & 0.773 [.746,.800] & 0.799 [.732,.857] & 0.875 [.737,.999] \\
PaDiM (in-dom.)      & 0.753 [.721,.783] & 0.761 [.699,.815] & 0.861 [.724,.992] \\
SumThreshold (flags) & 0.742 [.716,.767] & 0.911 [.861,.952] & 0.916 [.850,1.000] \\
SumThreshold (max)   & 0.811 [.786,.837] & 0.945 [.910,.972] & 0.966 [.944,1.000] \\
max-z fusion         & 0.839 [.822,.856] & 0.998 [.994,1.000] & 1.000 [.999,1.000] \\
\bottomrule
\end{tabular}
\end{table}

\begin{figure}[!t]
\centering
\includegraphics[width=\columnwidth]{injection_maps.png}
\caption{Injection--recovery maps on LOFAR (bottom row) and RSP1 (top row):
AUROC of injected-vs-base patches over RFI width (rows) and amplitude
(columns).  Down-arrows mark \emph{anti-correlated} cells
(AUROC $<0.5$): for the max-normalised SumThreshold statistic a wide,
moderate line raises its own per-line scale and \emph{suppresses} the
statistic elsewhere, ranking contaminated patches below clean ones.  MRD and both fusions darken from the weak-narrow corner; the champion
Fisher3$'$ (second column) additionally covers wide and transient
morphologies that MRD alone misses (cf.\ $w\ge3$ at $+20$--$50\%$).}
\label{fig:maps}
\end{figure}

\begin{figure}[!t]
\centering
\includegraphics[width=\columnwidth]{fig_roc_pdpfa.png}
\caption{Detection probability versus false alarm rate at patch level
(log $P_{fa}$ axis), full test set (left) and the persistent-narrowband
stratum (right).  At the operating points $P_{fa}=10^{-2}$/$10^{-3}$ the
max-z fusion attains $P_d=0.469$/$0.315$ on the full set, and on the
narrowband stratum the median-based detectors separate cleanly from all
area-based families throughout the low-false-alarm region.}
\label{fig:pdpfa}
\end{figure}

Patch-level $P_d$--$P_{fa}$ curves (Fig.~\ref{fig:pdpfa}) give the
operating-point view: at $P_{fa}=10^{-2}$ the fusion reaches $P_d=0.47$
(all patches) and the median-based detectors own the low-$P_{fa}$ region
of the narrowband panel outright.

\subsection{Injection maps: where each family breaks}
Fig.~\ref{fig:maps} summarises the width$\times$amplitude grids.  Four
patterns, each with a mechanism: (i) MRD and the fusion are
width-insensitive --- temporal aggregation asks whether a column persists,
not how wide it is --- and cross AUROC $0.7$--$0.9$ at $+20\%$ for every
width; (ii) the autoencoder and the industrial detectors respond to
\emph{area}: one-pixel lines are invisible at every amplitude on RSP1
(AUROC $0.500$), and only wide, strong injections register; (iii) the
classical flag count needs width but not strength ($w\ge2$, $+50\%$:
$0.8$--$1.0$; $w=1$: never); (iv) the max-normalised classical statistic is
\emph{anti-correlated} for wide moderate lines ($0.36$--$0.43$ at $w\ge3$,
$+20\%$) --- to our knowledge the first quantified instance of this trap.

\subsection{Absorption-like and multi-line injections}
Negative injections (channel power reduced by 10--50\%) show that MRD and
the fusion detect absorption-like RFI with the same sensitivity as emission
($-20\%$: $0.80$--$0.92$), while the autoencoder is blind across the entire
negative domain ($0.500$ even at $-50\%$).  With $n$ simultaneous one-pixel
lines the fusion is insensitive to $n$ ($0.90$ at $n=2$, unchanged at
$n=8$), the autoencoder remains at chance even for eight simultaneous lines,
and the classical flag count climbs with total area ($0.77\to1.00$) --- the
third independent confirmation of the complementary-coverage picture.

\subsection{Ablations}
\label{sec:ablation}
\textbf{Fusion members} (selection half): max-z two-member $0.834$;
Fisher three-member with the plain autoencoder $0.865$; with the denoising
autoencoder (DAE) $0.875$; adding the max-normalised SumThreshold statistic
\emph{lowers} it to $0.857$.  \textbf{DAE vs AE}: synthetic-RFI denoising
training alone lifts the autoencoder from $0.767$ to $0.794$ overall
(transient $+0.030$, wide $+0.046$, narrow $-0.015$).  \textbf{Skip-connection
AE}: trained on all 9217 clean patches it is better on transients but worse
on narrowband and does not improve the fusion --- reported as a negative
result.  \textbf{MRD scale set}: narrowband AUROC is $0.993$--$0.998$ for
every reasonable scale set; the multi-scale maximum contributes to the other
strata (overall $0.694\to0.724$).  \textbf{Calibration size}: the zero-label
protocol saturates with $\sim$100 clean patches (overall within $0.002$ of
the 4000-patch calibration) --- the data hunger of the method is negligible
for deployment.  \textbf{Network-engineering knobs}: a ten-variant grid on
the denoising member (losses L1/Huber/MSE+gradient/MSE+correlation;
AdamW and $10\times$ lower learning rate; cosine schedules; doubled width;
a residual refinement block at the bottleneck), each trained on identical
data order and identical synthetic-RFI draws, moves the member's overall
AUROC by at most $+0.001$ and $-0.065$ (L1) around the baseline, and the
best variant inside the full fusion reproduces the confirmation result
within bootstrap error ($0.912$ [0.896, 0.928] vs $0.914$ [0.898, 0.929]).
The member is insensitive to standard engineering knobs --- consistent with
the capacity, skip-connection, and member-search findings above: the
bottleneck is the aggregation paradigm, not the network.

\begin{figure}[!t]
\centering
\includegraphics[width=\columnwidth]{fig_forest.png}
\caption{Overall AUROC on the confirmation half (55 unseen images) with
image-level bootstrap 95\% CIs; single-member values (grey) are full-test
points for context.  The Fisher fusions separate from the strongest
classical baseline; the confirmed champion (Fisher4b) adds the official
PatchCore stack as a fourth member.}
\label{fig:forest}
\end{figure}

\subsection{Confirmation half and the optimised detector}
\begin{table}[!t]
\centering
\caption{Confirmation half (55 unseen LOFAR images, scored once): overall
AUROC and image-level 95\% CI.}
\label{tab:confirm}
\footnotesize
\setlength{\tabcolsep}{3pt}
\begin{tabular}{lcc}
\toprule
 & overall [CI95] & narrow \\
\midrule
ST-max (reference) & 0.824 [0.790, 0.854] & 0.974 \\
max-z fusion       & 0.859 [0.833, 0.881] & \textbf{1.000} \\
Fisher3 (MRD+AE+ST-flags) & 0.873 [0.849, 0.896] & 0.995 \\
Fisher3$'$ (MRD+DAE+ST-flags) & 0.889 [0.865, 0.909] & 0.996 \\
Fisher4b (+official PC) & \textbf{0.914} [0.898, 0.929] & 0.997 \\
\bottomrule
\end{tabular}
\end{table}
Table~\ref{tab:confirm} and Fig.~\ref{fig:forest}: the three-member
Fisher fusion separates from the strongest classical baseline
(0.889 [0.865, 0.909] vs 0.824 [0.790, 0.854]).  The narrowband-optimal
configuration remains the two-member max-z fusion (stratum $1.000$ on the
confirmation half); we report both and note the trade-off explicitly ---
Fisher combination buys overall and transient performance at a negligible
narrowband cost ($0.996$ vs $1.000$).  A fourth member closes the arc: once
The confirmation split is by dataset order (parent image = $i\,/\!/225$;
109 images split into first-half selection and second-half confirmation;
positive patches 4008 vs 4500), and on the full first-generation test set
the fusion's advantage over the strongest classical baseline is
significant under a paired per-image permutation test (mean
$\Delta$AUROC $+0.044$, $p=5\times10^{-4}$, $B{=}2000$).
A fourth member closes the arc: once
strongest single baseline (Section~\ref{sec:migration}), its own scores
enter the same clean-data Fisher combination --- the confirmation half
rises to $0.914$ [0.898, 0.929] with every stratum at or above $0.997$
except transient ($0.900$), and its lower CI bound clears the classical
baseline's upper bound.  The strongest-baseline-becomes-a-member pattern
now applies to both the classical (flag counts) and the industrial
(PatchCore) families. (The label \emph{max-z} denotes two related two-member variants across the selection-half and confirmation-half tables --- AE-max and AE-top16 respectively --- and the archived first-generation scores used for the paired test are not score-normalisation interchangeable with the re-run confirmation generation; both are released.)

\subsection{Runtime}
On 20 CPU threads, per $64\times64$ patch: ConvAE $0.37$\,ms, PaDiM $0.65$,
MRD $3.87$, max-z fusion $4.13$, PatchCore $4.79$, SumThreshold $31.18$.
The complete two-stage system is $7.5\times$ faster than the classical
pipeline while dominating it on the strata that matter for spectrometry.

\subsection{Third dataset (HERA): the boundary of the persistence assumption}
\label{sec:hera}
On the public HERA simulation corpus (280 labelled $512^2$ test images,
complex visibilities converted to power, RFI in 96\% of test patches ---
a fundamentally different regime from LOFAR's 33\%), the picture changes
qualitatively: no detector exceeds AUROC $0.68$ (tuned SumThreshold max
$0.676$; the fusion $0.649$; MRD $0.545$, narrow stratum $0.630$; all
members re-calibrated on HERA clean training data only).  The mechanism is
the mirror image of Section~V-B: when the interferer persists across the
\emph{entire} analysis window, the temporal median absorbs it into the
baseline and time aggregation carries no information --- only frequency-domain
contrast retains power.  This is not a failure of the method but the
quantitative boundary of its applicability: the coverage claim of this
paper holds for interferers that occupy \emph{less than the analysis
window} (the persistent stratum is defined at $>60\%$ occupancy); HERA's
simulation places most of its RFI outside that scope, and we report it as
such.  Two consistent observations remain: the denoising autoencoder
retrained in-domain beats the zero-shot transfer ($0.634$ vs $0.598$), and
no family dominates.

\subsection{Migration study: industrial anomaly detection}
\label{sec:migration}
The obvious modernisation --- treat RFI patch screening as industrial
anomaly detection --- was implemented and tested rather than assumed, at
three fidelity levels: our in-domain reimplementation (native $64{\times}64$,
overall $0.773$, narrow stratum $0.799$), the published ImageNet backbone at
native resolution ($0.666$; the domain gap), and finally the \emph{official}
anomalib PatchCore pipeline with its default preprocessing.  The official
stack upsamples every patch to $256{\times}256$ before the WideResNet-50
feature extractor --- and this changes the verdict: on a balanced
2000-patch subset (1000 RFI, 1000 clean; strata scored on identical
indices for all methods) it reaches overall $0.855$ with narrowband
$0.960$ (n=67), transient $0.838$, wide $0.936$.  A controlled decomposition falsifies our first guess at the
mechanism: with the same WRN50 backbone and the same
random-subsample memory bank, varying \emph{only} the input resolution
leaves the narrowband stratum unchanged ($0.779$ at $64$\,px native vs
$0.773$ at $256$\,px; overall $0.66\to0.71$).  The official stack's
narrowband capability comes from its \emph{fully-provisioned memory
bank} --- greedy coreset selection and neighbourhood reweighting --- not
from resolution.  The receptive-field-dilution blindness we report
therefore applies to simplified bank protocols (our reimplementation
arms), and the benchmark narrowband gap to our detectors narrows to
$0.96$ vs $0.99$ against the full official stack.

On the same balanced subset our detectors score: Fisher3$'$ $0.881$
(narrow $0.989$), max-z fusion $0.844$ ($0.999$), MRD $0.725$ ($0.999$),
ConvAE $0.790$ ($0.872$), SumThreshold max $0.808$ ($0.933$).  The
official implementation is thus genuinely competitive --- the strongest
single-model baseline in this paper --- and our advantage concentrates
exactly where the diagnosis predicts: sub-upsampling-contrast narrowband
($0.99$ vs $0.96$), zero-label deployment at native resolution, and CPU
runtime ($4.1$\,ms/patch at $64^2$ pixels versus a $256^2$ pipeline).

The \emph{most recent} published method replicates the same story.  RFI-HWT
(You et al., A\&A 701, A81, 2025 --- multilevel 2-D wavelet transform plus a
DnCNN denoiser per wavelet level) was executed from the authors' released
code and pretrained weights (three disclosed adaptations: the pulsar-FITS
wrapper bypassed, a released off-by-one in the inverse wavelet loop
restricted to the denoised levels, DataParallel prefix stripping; and a
fidelity note that $64{\times}64$ patches sit below the method's design
window, level-3 db8 coefficients being boundary-affected).  On the identical
balanced subset it is the strongest single-model baseline overall
($0.864$; narrow $0.877$) --- yet on the real instrument its one-pixel
injection ladder is chance at every contrast (AUROC $0.499$--$0.509$ at
$+10$--$100\%$), the third independent state-of-the-art method to confirm
that nothing bank- or denoiser-based survives the cross-instrument transfer
that the median-based stage handles with minutes of on-site quiet data.
One reproducibility note: the official loader assumes 8-bit imagery; our
first export in 16-bit saturated silently to AUROC $0.482$ --- a
cross-framework pitfall worth documenting.

\subsection{Real-instrument validation on RSP1}
\label{sec:rsp1val}
Applied to 600\,s RSP1 integrations, the MRD stage blindly discovers a
deterministic comb: 81 peaks at $z>5$ with best spacing $1086$ bins
$=66.28$\,kHz ($=2\times33.14$\,kHz, the Hann-FFT bin spacing of the
waterfall chain), autocorrelation $0.666$ ($10^{4}$-permutation and 64-channel-block-permutation
nulls both $p<1\mathrm{e}{-4}$, null p95 $0.113$/$0.114$); 22 of the 81 peaks sit
within $0.06$ of the comb lattice and carry $1.6\times$ the off-lattice mean
significance, and the comb is specific to the coherent chain --- absent in all
five independent incoherent-waterfall captures ($|{ac}|<0.006$ at the expected
lag).  Injected one-pixel lines are
ranked by the median-based stage with AUROC $0.57$ ($+5\%$) and $0.73$
($+10\%$), rising to $0.97$--$1.00$ at $+50$--$100\%$ after on-instrument
self-calibration of the fused members, while the
autoencoder, PatchCore, and PaDiM remain at chance at every amplitude and
the classical statistic transitions only at $+50\%$
(Fig.~\ref{fig:maps}, top row), and --- decisively --- the \emph{official}
anomalib PatchCore stack itself, deployed as-is with its LOFAR-trained
memory bank, scores AUROC $0.500$ on the same one-pixel ladder at every
amplitude from $+10\%$ to $+100\%$: the memory bank does not transfer
across instruments, whereas the median-based stage re-calibrates from
minutes of on-instrument quiet data (Section~V-G).  The instrument is thus both the stress test
and the demonstration that the pipeline finds real artefacts end to end.  The fused detector itself, calibrated on LOFAR-clean data and
applied cross-domain, inherits the weak-narrow corner and adds wide
coverage ($w{=}5$, $+20\%$: AUROC $0.94$ vs $0.83$ for MRD alone), but the
domain-shifted DAE member dilutes extreme narrowband ranking
($w{=}1$, $+100\%$: $0.85$ vs $0.96$); re-calibrating the members on the
instrument's own quiet data (249 patches, still label-free) recovers
$0.97$ and a $100\%$ detection rate at $+100\%$ while remaining at the
5\%-false-alarm operating point.  On-instrument self-calibration is
therefore the recommended deployment mode; adding the official-PatchCore
member (Fisher4b) leaves this ladder essentially unchanged
($0.72/0.82/0.84/0.85$ cross-domain and $0.64/0.82/0.89/0.97$ with
self-calibration at $+10$--$100\%$), confirming that the real-instrument
behaviour is carried by the median-based and denoising members, not by the
bank.

The newest zero-label paradigm --- promptable foundation-model segmentation
\cite{yang2025sam} --- fails the same transfer test.  On a fixed 64-patch
subsample of the identical injection ladder, the Segment Anything Model
(ViT-B) with automatic mask generation, scored blindly by the sharpest
per-channel spike among its output masks, ranks at chance to weak at every
contrast (AUROC $0.57/0.69/0.67$ at $+10/50/100\%$; $0\%$ detection at the
$5\%$-false-alarm operating point), while the median-based member on the
same patches scores $0.75/0.97/1.00$.  A point prompt placed on the injected channel anchors a dedicated line mask
only when the line is glaring: coverage at the line is $+1.00$ at $+100\%$
(dedicated strip in $32/32$ scans) but $+0.00$ at $+10\%$ --- and knowing
where to prompt presupposes the detection under test.  The blind path ---
the only zero-label deployment --- remains at chance at every amplitude.
Object priors without temporal persistence do not survive the
cross-instrument transfer in the mode that deployment requires; the oracle
mode is an answer-key check, not a detector.  The failure is
representational, not informational: re-rendering the aggregate-first
column-$z$ profile as a bar image lets the same model isolate the injected
line blindly (narrow-mask yield $0.88$--$0.95$; readout AUROC
$0.63/0.80/0.83$ at $+10/50/100\%$, up from $0.57/0.69/0.67$ on raw
patches) --- yet the readout remains strictly below the max-$|z|$ ceiling
of the very profile it consumes ($0.73/0.96/0.99$) at every amplitude.
Aggregation owns the information; the foundation model can only visualise
it.

\textbf{Transfer to an independent educational instrument.}  The suite also
transfers, zero-retuned, to a third-party educational horn-SDR corpus: 11
single-sweep galactic scans following the two-load protocol of Mhaske et
al.~\cite{mhaske2022}.  Single sweeps are the documented degenerate class
for per-scan statistics --- the spectra are signal-dominated smooth
products whose monotone runs zero the median-of-three noise estimator ---
and the survey-standard fix, cross-scan common-mode rejection (the
per-channel median across the 11 scans), resolves it:  the corpus's known
RFI birdie ($-168.6$\,km/s, $l=5$ scan) is the largest scan-unique feature
($10.7\sigma$; leave-one-out $10.8\sigma$), while the instrumental
structure that dominates per-scan statistics ($17$--$40\sigma$) collapses
below $4\sigma$ by construction.  Aggregation owns the detection here too
--- across scans rather than across time.

\section{Discussion}
\subsection{Improvements over the replicated baselines}
\label{sec:improve}
The three baseline families were not taken from leaderboards but
\emph{replicated}: SumThreshold from its published equations
\cite{offringa2010}, the RFI-NLN autoencoder pipeline executed end-to-end
from the authors' original code on the same public pickle (reduced scale:
750 of 7465 training images, 10 of 100 epochs, CPU-only; pixel AUROC
$0.771$ vs their published $0.840$ at full scale, with the deviation
disclosed), and PatchCore/PaDiM from their published recipes including the
original backbone.  Replication at this level is what licenses the
improvement claims below: the baselines keep their published
hyperparameters and our gains are measured against what their methods
actually do, not against tuned variants.

\emph{Versus RFI-NLN.}  The replicated pipeline is faithful, and its
blind spot is structural: a strided encoder averages a one-pixel line away
at the first level, so no amount of latent-neighbour refinement (their NLN
scoring) can recover a feature the encoder never represented --- on our
injection grid the AE family sits at AUROC $0.500$ for one-pixel lines at
every contrast.  Three changes, all label-free: (i) a temporal-median
member that aggregates along time \emph{before} any pooling, restoring
exactly the persistence information the encoder destroys (narrowband
stratum $0.998$ vs $0.869$); (ii) the deep member retrained as a denoiser
with synthetic RFI of matched morphology rather than as a pure
reconstruction autoencoder ($+0.027$ overall on the same benchmark); and
(iii) top-$k$ rather than max patch aggregation, which buys area-robustness
($+0.015$).  The reduced-scale replication also showed the direction is
robust: even at $1/10$ of the training budget the AE family never
approaches the fusion on the narrow stratum.

\emph{Versus SumThreshold.}  We changed the algorithm not at all --- we
demoted it.  As a standalone detector its shared threshold ladder floors
persistent sub-threshold lines (chance-level ranking below $+50\%$
contrast) and its max-normalised statistic \emph{anti-correlates} for wide
moderate lines ($0.36$--$0.43$).  Inside the Fisher fusion, however, its
flag count contributes exactly the area sensitivity that the median-based
members lack ($w\ge2$ lines at $+50\%$: $0.8$--$1.0$ where MRD alone
transfers less), and the clean-data p-value transform absorbs its
per-member miscalibration.  The lesson is architectural: the classical
workhorse is most valuable as a \emph{member}, not as the detector.

\emph{Versus PatchCore/PaDiM.}  With the published backbone the
ImageNet$\to$spectrogram domain gap costs $0.11$ AUROC; with in-domain
features the gap closes but one-pixel injections remain at chance on the
real instrument --- receptive fields dilute what a $32\times32$ crop of
columns concentrates \emph{under simplified bank protocols}; the
fully-provisioned official stack recovers most of the benchmark gap
(Section~\ref{sec:migration}) yet does not transfer to the real
instrument at any contrast --- exactly where temporal aggregation,
which needs no bank at all, wins.  The improvement here is not a better memory bank
but an earlier stage: temporal aggregation \emph{before} patch embedding.
Equivalently: the memory-bank family should be applied to
time-aggregated residuals, not to raw spectrograms --- a prescription that
follows directly from the diagnosis and is testable on any industrial-AD
stack.

The common denominator of all three replications is that every member
computes statistics without aggregation along the time axis.  The two-stage
design is therefore not one more family but a re-ordering: aggregate first,
then apply each family where its inductive bias matches the morphology.

\subsection{One mechanism, five families}
Every failure documented here is a corollary of a single statement: a
statistic that does not aggregate along the time axis carries no information
about persistence, and a persistent line below the multiple-testing floor is
indistinguishable from noise per sample.  The deep family fails through
stride averaging, the classical family through its shared threshold ladder,
the industrial family through receptive-field dilution, the foundation-model
family through object priors that carry no temporal persistence
(Section~\ref{sec:rsp1val}) --- and a 1960s-vintage
time median succeeds.  For spectrometry, where narrowband carriers are the
dominant morphology, this ordering \emph{reverses} the usual deep-versus-
classical narrative.

\subsection{A floor theorem for threshold ladders}
\label{sec:floor}
The mechanism can be stated exactly.  Let $R(\nu,t)$, $\nu{=}1..F$,
$t{=}1..T$, be i.i.d. $\mathcal{N}(0,\sigma^2)$ samples and let one channel
$\nu_0$ carry a persistent line of amplitude $\delta$,
$R(\nu_0,t)=\delta+\text{noise}$.  A threshold-ladder detector flags windows
of size $M$ along $t$ when the window mean exceeds
$\chi_1\,\rho^{-\log_2 M}\sigma$ ($\rho>1$ fixed, Offringa's Eq.~12), where
$\chi_1$ is calibrated for a patch-level false alarm $\alpha$ over the
$F\cdot T$ single-sample tests.

\emph{Proposition 1 (ladder floor).}  (i) $\chi_1 \ge
\Phi^{-1}(1-\alpha/2FT) \sim \sigma\sqrt{2\ln(FT/\alpha)}$.  (ii) For a
persistent line the optimal ladder window is $M{=}T$ (window means all
estimate $\delta$ regardless of $M$, while $\chi_M$ decreases with $M$), so
the ladder detects iff
\begin{equation}
\delta \;\ge\; \delta^{\mathrm{ST}}_{\min}
 \;=\; \chi_1\,\sigma\,\rho^{-\log_2 T}
 \;\sim\; \sigma\sqrt{2\ln(FT/\alpha)}\;T^{-\log_2\rho}.
\label{eq:floor}
\end{equation}
(iii) The time-median statistic of the line channel is
$\delta + O(\sigma/\sqrt{T})$ (Gaussian efficiency $1.2533$), and the
$\max$ over $F$ null channels is $\sim\sigma\sqrt{2\ln F}/\sqrt{T}$, so the
median detects iff
\begin{equation}
\delta \;\ge\; \delta^{\mathrm{MRD}}_{\min}
 \;\sim\; \sigma\sqrt{2\ln F / T}.
\label{eq:mrd}
\end{equation}
(iv) The ratio $\delta^{\mathrm{ST}}_{\min}/\delta^{\mathrm{MRD}}_{\min}
\sim \sqrt{T}\,T^{\log_2\rho}\cdot\mathrm{const} \to \infty$: the ladder's
false-alarm floor, set at $M{=}1$, consumes exactly the $\sqrt{T}$
integration that the median realises.  The same floor applies to
\emph{any} per-sample or per-position statistic --- including strided
network features --- because none of them escapes the $F\cdot T$
multiple-testing burden; temporal aggregation is the only escape.

\emph{Proof sketch.}  (i) Union bound over $FT$ Gaussian tests.  (ii) Every
window mean of the line channel estimates $\delta$ with noise
$\sigma/\sqrt{M} \le \sigma$, so the binding constraint is the threshold,
minimised at $M{=}T$.  (iii) Median concentration plus the union bound over
$F$ channels.  (iv) Direct quotient of
\eqref{eq:floor} and \eqref{eq:mrd}.  \hfill$\square$

\emph{Numerical instantiation.}  With $T{=}F{=}64$, $\rho{=}1.5$,
$\alpha{=}10^{-3}$, and the measured per-sample log-noise
$\sigma\!\approx\!0.5$: Proposition~1 predicts
$\delta^{\mathrm{ST}}_{\min}\!\approx\!12.8/11.4\!\approx\!1.12\,\sigma$
--- i.e., the classical ladder turns on between $+50\%$ ($0.81\sigma$) and
$100\%$ ($1.38\sigma$) --- and
$\delta^{\mathrm{MRD}}_{\min}\!approx\!0.21\sigma$ --- with
$F_{\mathrm{eff}}\!approx\!4$, the number of independent sign-runs of the
double-median residual over a 64-channel window (the residual vanishes
identically on monotone runs; a naive $F{=}64$ over-counts and would give
$0.36\sigma$) --- turning on below
$10\%$ ($0.14\sigma$): exactly the measured injection turning points of
Fig.~\ref{fig:maps}.  The theorem also explains why the DAE and
PatchCore-style members remain at chance on the same injections
(Section~\ref{sec:migration}): stride and receptive-field statistics are
per-position statistics and inherit the same $F\cdot T$ floor, whereas the
temporal-median member is the only one whose effective sample size grows
A three-working-point scan of the threshold multiplier ($\chi_1\times0.5/1/2$ around the calibrated value) confirms the floor's $\chi_1$ dependence empirically: one-pixel AUROC at $+100\%$ degrades monotonically $0.93/0.86/0.65$, and no working point turns on below $+50\%$ --- the deployed $\chi_1$ is the calibrated one, not a weakened one.
with $T$.

\subsection{Design guidance for low-cost spectrometers}
(i) Aggregate along time before anything else; (ii) if a deep stage is
desired, train it as a denoiser with synthetic morphology --- it costs
minutes on CPU and needs no labels; (iii) combine members by clean-data
Fisher p-values rather than score maxima; (iv) never trust a single
max-normalised statistic for wide weak interferers (the anti-correlation
trap); (v) validate with injection maps, not with a single pooled number.

\subsection{Limitations}
The SumThreshold baseline implements the plain published variant without the
authors' flag-exclusion refinement; the HERA corpus bounds the method's
applicability to sub-window persistent interferers
(Section~\ref{sec:hera}); RSP1 provides one real instrument, single
polarisation; the DAE's synthetic morphology is defined by the authors and
could be expanded; the foundation-model arm is single-scale (ViT-B) and
real-instrument-only; and the anti-multiplicity discipline relies on a single
pre-registered confirmation split.

\section{Conclusion}
Weak persistent narrowband interference --- the morphology that matters most
for hydrogen-line spectrometry --- is invisible to deep, classical, and
industrial detector families alike, each for a different structural reason
that injection maps make directly measurable.  A label-free
Fisher fusion whose members own the disjoint blind regions recovers it:
with the official PatchCore stack as a calibrated fourth member,
$0.997$ AUROC on persistent narrowband, overall $0.914$ [0.898, 0.929] on
held-out images
with CI separation from the strongest classical baseline, blind comb
discovery and injection recovery on a real EUR~200 instrument, at
$7.5\times$ the speed of the classical pipeline.  The protocol --- stratified
injection maps, cluster-bootstrap CIs, selection/confirmation discipline ---
is instrument-agnostic and released for reuse.

\section*{Acknowledgment}
% TODO[1]
The author thanks the RFI-NLN authors for the public LOFAR dataset and
labelled masks.

\section*{Code and Data Availability}
All experiments run on public data: the LOFAR and HERA benchmark
(Zenodo record 8275061), the LAB survey cube, and the community
corpora listed in the paper, each with their original licences.  The
detector suite (MRD, denoising AE, SumThreshold reference
implementation, PatchCore/PaDiM arms, fusion, evaluation protocol),
the reduced-scale runner for the original RFI-NLN code (with the four
disclosed compatibility patches), the anomalib export harness, and
every result file behind each table and figure (JSON with
machine-readable provenance) are archived at
\texttt{[repository DOI upon acceptance]}.  The RSP1 raw IQ files are
available from the author on reasonable request (instrument-owner
data).

\begin{thebibliography}{99}
\bibitem{offringa2010} A.~R. Offringa \emph{et al.}, ``Post-correlation
radio frequency interference classification methods,'' \emph{MNRAS},
vol.~405, no.~1, pp.~155--167, 2010, doi:10.1111/j.1365-2966.2010.16471.x.
\bibitem{offringa2012} A.~R. Offringa \emph{et al.}, ``The LOFAR radio
environment,'' \emph{A\&A}, vol.~549, A11, 2012, doi:10.1051/0004-6361/201220293.
\bibitem{offringa2013} A.~R. Offringa, ``The SumThreshold method: technical
details,'' technical report, 2013. [Online]. Available:
\url{https://andreoffringa.org/pdfs/SumThreshold-technical-report.pdf}
\bibitem{akeret2017} J.~Akeret \emph{et al.}, ``Radio frequency interference
mitigation using deep convolutional neural networks,'' \emph{Astronomy and
Computing}, vol.~18, pp.~35--39, 2017, arXiv:1609.09077.
\bibitem{vafaieisadr2020} S.~Vafaei Sadr \emph{et al.}, ``Deep learning
improves identification of radio frequency interference,'' \emph{MNRAS},
vol.~499, pp.~379--390, 2020, doi:10.1093/mnras/staa2724.
\bibitem{mesarcik2022} M.~Mesarcik \emph{et al.}, ``Learning to detect
radio frequency interference in radio astronomy without seeing it,''
\emph{MNRAS}, vol.~516, pp.~5367--5378, 2022, doi:10.1093/mnras/stac2503.
\bibitem{mesarcik2023} M.~Mesarcik \emph{et al.}, ``The ROAD to discovery:
machine-learning-driven anomaly detection in radio astronomy
spectrograms,'' \emph{A\&A}, vol.~680, A74, 2023,
doi:10.1051/0004-6361/202347182.
\bibitem{dutoit2024} C.~Du Toit \emph{et al.}, ``A comparison framework for
deep-learning-based RFI detection,'' \emph{MNRAS}, vol.~530, pp.~613--627,
2024.
\bibitem{pritchard2024} D.~Pritchard, A.~Wicenec, M.~Bennamoun, and
S.~Dodson, ``RFI detection with spiking neural networks,'' \emph{PASA},
vol.~41, e028, 2024, doi:10.1017/pasa.2024.27.
\bibitem{rfihwt2025} Z.-Y. You \emph{et al.}, ``RFI-HWT: hybrid deep-learning
algorithm for advanced RFI mitigation in radio astronomy,'' \emph{A\&A},
vol.~701, A81, 2025.
\bibitem{he2025} J.~He, Y.~Li, and X.~Chen, ``A review of software-based
detection and mitigation techniques for radio frequency interference,''
\emph{Astronomical Techniques and Instruments}, vol.~3, pp.~312--336, 2025,
doi:10.3724/ati2025048.
\bibitem{roth2022} K.~Roth, L.~Pemula, J.~Zepeda, B.~Sch\"olkopf,
T.~Brox, and P.~Gehler, ``Towards total recall in industrial anomaly
detection (PatchCore),'' in \emph{Proc.~IEEE/CVF CVPR}, 2022, pp.~14298--14308,
doi:10.1109/CVPR52688.2022.01392.
\bibitem{defard2021} T.~Defard, A.~Setkov, A.~Loesch, and R.~Audigier,
``PaDiM: a patch distribution modeling framework for anomaly detection
and localization,'' in \emph{Proc.~ICPR Int.\ Workshops (LNCS)}, 2021,
pp.~475--489, doi:10.1007/978-3-030-68799-1\_35.
\bibitem{saraos2026} C.~H.~A. Saraos, K.~M. Senados, and E.~N. Hamoy,
``Low-cost 21\,cm hydrogen-line detection for physics education using
software-defined radio and a simple antenna,'' \emph{Phys.\ Educ.},
vol.~61, 045025, 2026, doi:10.1088/1361-6552/ae806b.
\bibitem{pnu2025} A.~Allawi \emph{et al.}, ``Design and optimization of a
low-cost 5-m radio telescope at Princess Nourah University (PNU), Saudi
Arabia,'' \emph{Radio Science}, vol.~60, no.~2, e2024RS008170, 2025,
doi:10.1029/2024RS008170.
\bibitem{phelps2024} C.~Phelps, ``Galactic neutral hydrogen structures
spectroscopy and kinematics: designing a home radio telescope for 21\,cm
emission,'' arXiv:2411.00057, 2024 (unpublished preprint).
\bibitem{connor2019} E.~F. Keane and C.~R.~H. Walker, ``Fast radio burst
injection tests,'' arXiv:1907.07948, 2019.
\bibitem{haynes2018} M.~P. Haynes \emph{et al.}, ``The Arecibo Legacy Fast
ALFA survey: the ALFALFA extragalactic H\,I source catalog,''
\emph{ApJ}, vol.~861, no.~1, art.~49, 2018, doi:10.3847/1538-4357/aac956.
\bibitem{yang2025sam} Y.~Yang, F.~Zhao, R.~Liang, Q.~Guo, J.~Gu, Y.~Huang,
and Y.~Yu, ``Performance of the segment anything model in various
RFI/events detection in radio astronomy,'' \emph{Publ. Astron. Soc. Aust.},
vol.~42, e019, 2025, doi:10.1017/pasa.2024.128.
\bibitem{mhaske2022} A.~A. Mhaske, J.~Bagchi, B.~C. Joshi, J.~Jacob, and
P.~K. T., ``A Bose horn antenna radio telescope (BHARAT) design for 21 cm
hydrogen line experiments for radio astronomy teaching,''
\emph{Am. J. Phys.}, vol.~90, no.~12, pp.~948--956, 2022,
doi:10.1119/5.0065381.
\end{thebibliography}

\end{document}
